\SourceLeaf{325}%
there and no longer appear in the results, may be used here for a different purpose without difficulty.

Draw a geodesic line, initially of unspecified extent, through the two points on the surface of the ellipsoid to which the problem relates, and choose an arbitrary starting point on it. Denote the part of this line from the starting point to an unspecified point by \(u\). At the latter point, take the geodesic in the direction of increasing \(u\) and the meridian from north to south, and denote by \(X\) the angle the geodesic makes with the meridian. Let \(Y\) and \(Z\) denote the latitude and longitude of the unspecified point. I assume that longitudes increase from west to east and azimuths \(X\) from south toward west. If, as throughout, the semimajor axis and eccentricity of the generating ellipse are denoted by \(a\) and \(e\), we have, for well-known reasons,
\[
\begin{aligned}
\frac{\mathrm{d}Y}{\mathrm{d}u}&=-\frac{\cos X(1-ee\sin Y^2)^{\frac32}}{(1-ee)a}\quad .\ .\ .\ .\ .\ .\ .\ .\ .\ .\quad(1)\\
\frac{\mathrm{d}Z}{\mathrm{d}u}&=-\frac{\sin X(1-ee\sin Y^2)^{\frac12}}{a\cos Y}\quad .\ .\ .\ .\ .\ .\ .\ .\ .\ .\quad(2)
\end{aligned}
\]
\Needspace{8\baselineskip}
Furthermore, by a well-known theorem, the quantity
\[
\frac{\sin X\cos Y}{\BareRoot{(1-ee\sin Y^2)}}
\]
is constant for all points on the same geodesic line. Logarithmic differentiation thus gives
\[
\operatorname{cotang}X\,\mathrm{d}X=(\operatorname{tang}Y-\frac{ee\cos Y\sin Y}{1-ee\sin Y^2})\mathrm{d}Y=\frac{(1-ee)\operatorname{tang}Y}{1-ee\sin Y^2}\mathord{.}\mathrm{d}Y
\]
and hence, on combining this with (1),
\[
\frac{\mathrm{d}X}{\mathrm{d}u}=-\frac{\sin X\operatorname{tang}Y(1-ee\sin Y^2)^{\frac12}}{a}\quad .\ .\ .\ .\ .\ .\ .\ .\ .\quad(3)
\]

\indent We will, however, take our problem in a more general form and set
\[
\frac{\mathrm{d}X}{\mathrm{d}u}=x,\qquad\frac{\mathrm{d}Y}{\mathrm{d}u}=y,\qquad\frac{\mathrm{d}Z}{\mathrm{d}u}=z
\]
assuming for now only that \(x,y,z\) are arbitrary given functions of the two variables \(X\) and \(Y\). Further differentiation gives
